\chapter{Реализация и тестване}

\section{Реализация на класове}
В основата на реализацията на електронната таблица стои полиморфният контейнер (правоъгълна матрица от клетки), управляван от класа \texttt{Spreadsheet}. Тъй като таблицата трябва да съхранява хетерогенни данни (текст, числа и формули) в обща структура, тя е реализирана като двумерен вектор от указатели към базовия абстрактен клас \texttt{Cell}:

\begin{lstlisting}[language=C++]
// Декларация на основния контейнер в класа Spreadsheet
std::vector<std::vector<Cell *>> cells;
\end{lstlisting}

Използването на указатели към базовия клас позволява прилагането на полиморфизъм. Когато \texttt{Spreadsheet} изпълнява команди като отпечатване или изчисляване, той извиква виртуалните методи (като \texttt{print()} или \texttt{getNumericValue()}) директно през общия интерфейс на \texttt{Cell}. По този начин главният клас управлява всички видове клетки по абсолютно еднакъв начин, без нуждата да проверява ръчно дали конкретният обект зад указателя е число, текст или формула.

За да се попълни тази матрица с реални обекти при десериализация или редактиране, е имплементиран частният метод \texttt{parseCell}. Тъй като останалата част от логиката работи само с абстрактни указатели, този метод служи като мост – той директно приема текстовия низ, разпознава типа му чрез помощни функции и заделя памет за съответния наследник:

\begin{lstlisting}[language=C++]
// Метод за разпознаване и създаване
// на конкретни обекти клетки
Cell* parseCell(const std::string &token, Address a);
\end{lstlisting}

По този начин процесът по разпознаване на типовете е капсулиран на едно единствено място в кода, което улеснява бъдещото добавяне на нови видове клетки.

\subsection*{Реализация на връзката между таблицата и формулите}
Едно от основните имплементационни предизвикателства е взаимодействието между класа \texttt{Spreadsheet} и обектите от тип \texttt{Formula}. Първоначалната идея за осигуряване на достъп на формулите до клетките беше подаване на директен указател към таблицата, за да се предотврати копирането на целия двумерен вектор. Тъй като обаче формулите нямат реална нужда от достъп до всички клетки, бе реализирано по-елегантно решение чрез \texttt{std::function}.

Основният механизъм, който прави това възможно, е т.нар. \textit{capture list} (списък за улавяне) в C++. Чрез него ламбда изразът в \texttt{solveFormula} „улавя“ указателя към текущата инстанция на таблицата (\texttt{[this]}) и го предава индиректно:

\begin{lstlisting}[language=C++]
// Ламбда функция за индиректно извличане на стойности
formula.solveFormula(tokens, operators, [this](Address a) {
    return this->getOrCalculateCellValue(a);
});
\end{lstlisting}

По този начин обектът \texttt{Formula} извлича стойности от адреси в реално време, без изобщо да се интересува от структурата на двумерния масив и без да нарушава капсулацията на \texttt{Spreadsheet}.

\subsection*{Използван външен код}
Логиката във функцията \texttt{TextCell::getContentLength()} във файла \texttt{cell.cpp} е взета от генеративен модел, тъй като стандартният метод \texttt{std::string::length()} не работи коректно за текст на кирилица.

\section{Управление на паметта, алгоритми и оптимизации}

\textbf{Управление на ресурсите чрез „Голямата четворка“} \\
Тъй като \texttt{Spreadsheet} управлява вектор от указатели към динамично заделена памет, в него е реализирана „Голямата четворка“ (\textit{Big Four}) -- дефинирани са конструктор по подразбиране, копиращ конструктор, оператор за присвояване и деструктор. За нуждите на копирането се използва виртуалният метод \texttt{clone()} от \texttt{Cell}. Всеки наследник връща свое дълбоко копие (\textit{deep copy}), което елиминира риска от споделена памет и предпазва от „двойно освобождаване“ (\textit{double-free}).

\textbf{Полиморфно разрушаване на обекти} \\
Задължителен елемент от йерархията е дефинирането на виртуален деструктор в базовия клас \texttt{Cell}. Той гарантира, че при разрушаване на клетка през указател от базов тип, правилно ще се извърши полиморфно извикване на деструктора на конкретния наследник, предотвратявайки изтичане на памет (\textit{memory leak}).

\textbf{Изчисляване на формули и защита от безкрайна рекурсия} \\
Изчисляването на формули с препратки към други клетки изисква таблицата да бъде предварително заредена в паметта, тъй като даден израз може да реферира към клетка, която все още не е пресметната. Поради този факт стандартното линейно изчисляване по време на самото четене е неприложимо. Процесът е разделен на два етапа: първо се зареждат всички независими клетки (числа, текст и формули без препратки), след което се прилага стратегия за отложено изчисляване -- при срещане на препратка във формула, алгоритъмът се извиква рекурсивно за съответната клетка, за да извлече стойността ѝ в реално време.

За да се предотврати безкрайна рекурсия при циклични зависимости (напр. две клетки, които взаимно се реферират), преди всяко рекурсивно извикване текущата клетка се замества с \texttt{EmptyCell} и се маркира временно със статус \texttt{"CALCULATING"} (това е конкретното предназначение на член-данната \texttt{status} в \texttt{EmptyCell}). При повторен опит за достъп до клетка в това състояние се индикира зацикляне и изпълнението се прекъсва безопасно чрез хвърляне на изключение.

\textbf{Оптимизация на матрицата} \\
При четене от файл редовете често съдържат различен брой клетки, което прави структурата в паметта „назъбена“ (\textit{jagged}). Методът \texttt{alignAllCells} допълва матрицата с празни клетки (\texttt{EmptyCell}) до изравняване на всички редове спрямо най-дългия. Това осигурява правилна правоъгълна структура и елиминира риска от грешки при индексиране извън границите.

\section{Планиране и тестови сценарии}

Коректността на системата е проверена чрез примерната таблица, показана на табл. \ref{tab:test_table}.

\begin{table}[h!]
  \centering
  \begin{tabular}{|c|c|c|c|}
    \hline
    10 & 20 & 30 & 40 \\ \hline
    &    &    &    \\ \hline
    10 &    & 1000 &  \\ \hline
    &    &    &    \\ \hline
    & 10 &    &    \\ \hline
  \end{tabular}
  \caption{Начално състояние на тестовата таблица}
  \label{tab:test_table}
\end{table}

Върху нея през конзолния интерфейс са изпълнени сценарии, покриващи нормалната работа и защитите от грешки:

\begin{itemize}
  \item \textbf{Сценарий 1: Изчисляване на формула с валидни адреси} \\
    \textit{Вход:} \texttt{edit R2C1 =R1C1+R1C2} \\
    \textit{Резултат:} Системата успешно събира стойностите (10 и 20). При отпечатване клетката извежда \texttt{30}.

  \item \textbf{Сценарий 2: Операция с празна клетка} \\
    \textit{Вход:} \texttt{edit R2C2 =R3C1+R3C2} \\
    \textit{Резултат:} Тъй като \texttt{R3C2} е празна (интерпретира се като 0), резултатът от израза се оценява коректно на \texttt{10}.

  \item \textbf{Сценарий 3: Защита от деление на нула} \\
    \textit{Вход:} \texttt{edit R4C1 =R1C4/R2C4} \\
    \textit{Резултат:} Клетката \texttt{R2C4} е празна (0). Математическата валидация улавя делението на нула, предпазва програмата от срив и клетката връща \texttt{ERROR}.

  \item \textbf{Сценарий 4: Защита от циклична зависимост} \\
    \textit{Вход:} \texttt{edit R2C1 =R3C3+R2C1} \\
    \textit{Резултат:} Алгоритъмът засича статус \texttt{"CALCULATING"} заради самореференцията, прекъсва цикъла безопасно.
\end{itemize}